\SeventhArticle{16}
The whole task requires us to find, for any given value of $n$, a function $T$ of the form $t^{n+1}+\alpha t^n+\alpha't^{n-1}+\alpha''t^{n-2}$, etc., so constructed that, when the product
\[
T(t^{-1}+\tfrac12t^{-2}+\tfrac13t^{-3}+\tfrac14t^{-4}+\text{etc.})
\]
is expanded, the powers $t^{-1},t^{-2},t^{-3}\mathinner{.\,.\,.\,.}t^{-(n+1)}$ all have the coefficient $0$; or, if one prefers, a function $U$ of the form $u^{n+1}+\beta u^n+\beta'u^{n-1}+\beta''u^{n-2}+$, etc., whose product with $u^{-1}+\frac13u^{-3}+\frac15u^{-5}+\frac17u^{-7}+$, etc., is free of the powers \SourceLeaf{193}%
$u^{-1},u^{-2},u^{-3},u^{-4}\mathinner{.\,.\,.\,.}u^{-(n+1)}$. The latter method will be somewhat simpler: for it is readily seen that, to satisfy the prescribed condition, the coefficients of $U$ must vanish alternately, or that $\beta=0,\beta''=0,\beta''''=0$, etc., must be set; hence nearly half the work may be considered done already. Let us work out some of the simpler cases.

I. For $n=0$, the single coefficient of $t^{-1}$ in the product
\[
(t+\alpha)(t^{-1}+\tfrac12t^{-2}+\tfrac13t^{-3}+\text{etc.})
\]
must vanish. Since it becomes $=\frac12+\alpha$, we have $\alpha=-\frac12$, or $T=t-\frac12$. Likewise, $U=u$.

II. For $n=1$, the determination of $T$ depends on two equations
\[
\begin{aligned}
0&=\tfrac13+\tfrac12\alpha+\alpha'\\
0&=\tfrac14+\tfrac13\alpha+\tfrac12\alpha'
\end{aligned}
\]
from which we deduce $\alpha=-1,\alpha'=+\frac16$, or $T=tt-t+\frac16$. The determination of the function $U$ gives just one equation
\[
0=\tfrac13+\beta'
\]
from which $\beta'=-\frac13$, or $U=uu-\frac13$.

III. For $n=2$, the function $T$ is determined with the aid of three equations
\[
\begin{aligned}
0&=\tfrac14+\tfrac13\alpha+\tfrac12\alpha'+\alpha''\\
0&=\tfrac15+\tfrac14\alpha+\tfrac13\alpha'+\tfrac12\alpha''\\
0&=\tfrac16+\tfrac15\alpha+\tfrac14\alpha'+\tfrac13\alpha''
\end{aligned}
\]
from which we obtain $\alpha=-\frac32,\alpha'=\frac35,\alpha''=-\frac1{20}$, and therefore $T=t^3-\frac32tt+\frac35t-\frac1{20}$. To determine $U$, a single equation suffices
\[
0=\tfrac15+\tfrac13\beta'
\]
from which $\beta'=-\frac35$, or $U=u^3-\frac35u$.

However, we shall not pursue this method further here, since it leads to increasingly troublesome calculations, but shall proceed to the true source of the general solution.
\SourceLeaf{194}
\SeventhArticle{17}
Given a continued fraction
\[
\varphi=\cfrac{v}{w+\cfrac{v'}{w'+\cfrac{v''}{w''+\cfrac{v'''}{w'''+\text{etc.}}}}}
\]
it is known that successively closer approximating fractions are found by the following algorithm. Form two series of quantities, $V,V',V'',V'''$, etc., $W,W',W'',W'''$, etc., by these formulas
\[
\begin{array}{ll}
V=0&W=1\\[4pt]
V'=v&W'=wW\\[4pt]
V''=w'V'+v'V&W''=w'W'+v'W\\[4pt]
V'''=w''V''+v''V'&W'''=w''W''+v''W'\\[4pt]
V''''=w'''V'''+v'''V''&W''''=w'''W'''+v'''W''
\end{array}
\]
etc.; then we shall have
\[
\begin{aligned}
\cfrac VW&=0\\[9pt]
\cfrac{V'}{W'}&=\cfrac vw\\[9pt]
\cfrac{V''}{W''}&=\cfrac v{w+\cfrac{v'}{w'}}\\[9pt]
\cfrac{V'''}{W'''}&=\cfrac v{w+\cfrac{v'}{w'+\cfrac{v''}{w''}}}
\end{aligned}
\]
and so on. Further, it is known, or is easily confirmed from the preceding equations themselves, that
\[
\begin{aligned}
VW'-V'W&=-v\\
V'W''-V''W'&=+vv'\\
V''W'''-V'''W''&=-vv'v''\\
V'''W''''-V''''W'''&=+vv'v''v'''
\end{aligned}
\]
etc. Hence the following holds for the series
\[
\frac v{WW'}-\frac{vv'}{W'W''}+\frac{vv'v''}{W''W'''}-\frac{vv'v''v'''}{W'''W''''}+\text{etc.}
\]
\SourceLeaf{195}
\[
\begin{array}{ll}
\text{first term}&=\dfrac{V'}{W'}\\[10pt]
\text{sum of the first two terms}&=\dfrac{V''}{W''}\\[10pt]
\text{sum of the first three terms}&=\dfrac{V'''}{W'''}\\[10pt]
\text{sum of the first four terms}&=\dfrac{V''''}{W''''}
\end{array}
\]
and so on; therefore the series itself, continued either to infinity or until it terminates, will express the continued fraction $\varphi$ itself. This also gives the difference between $\varphi$ and each of the approximating fractions $\frac{V'}{W'},\frac{V''}{W''},\frac{V'''}{W'''}$, etc.

From formula 33 of Article 14 of the \textit{General Investigations of the Infinite Series}, by changing $t$ to $\frac1u$, we easily obtain the transformation of the series
\[
\varphi=u^{-1}+\tfrac13u^{-3}+\tfrac15u^{-5}+\tfrac17u^{-7}+\text{etc.}
\]
into the following continued fraction
\[
\cfrac{1}{u-\cfrac{\cfrac13}{u-\cfrac{\cfrac{2\mathbin{.}2}{3\mathbin{.}5}}{u-\cfrac{\cfrac{3\mathbin{.}3}{5\mathbin{.}7}}{u-\cfrac{\cfrac{4\mathbin{.}4}{7\mathbin{.}9}}{u-\text{etc.}}}}}}
\]
so that we have
\[
\begin{aligned}
v&=1,\ v'=-\tfrac13,\ v''=-\tfrac4{15},\ v'''=-\tfrac9{35},\ v''''=-\tfrac{16}{63}\ \text{etc.}\\
w&=w'=w''=w'''=w''''\ \text{etc.}=u.
\end{aligned}
\]
Hence for $V,V',V'',V'''$, etc., $W,W',W'',W'''$, etc., we obtain the following values
\[
\begin{aligned}&V=0,\\[5pt]&\qquad W=1\\[5pt]&V'=1,\\[5pt]&\qquad W'=u\\[5pt]&V''=u,\\[5pt]&\qquad W''=uu-\tfrac13\\[5pt]&V'''=uu-\tfrac4{15},\\[5pt]&\qquad W'''=u^3-\tfrac35u\end{aligned}\]
\[\begin{aligned}&V''''=u^3-\tfrac{11}{21}u,\\[5pt]&\qquad W''''=u^4-\tfrac67uu+\tfrac3{35}\\[5pt]&V^{\mathrm{V}}=u^4-\tfrac79uu+\tfrac{64}{945},\\[5pt]&\qquad W^{\mathrm{V}}=u^5-\tfrac{10}{9}u^3+\tfrac5{21}u\\[5pt]&V^{\mathrm{VI}}=u^5-\tfrac{34}{33}u^3+\tfrac15u,\\[5pt]&\qquad W^{\mathrm{VI}}=u^6-\tfrac{15}{11}u^4+\tfrac5{11}uu-\tfrac5{231}\\[5pt]&V^{\mathrm{VII}}=u^6-\tfrac{50}{39}u^4+\tfrac{283}{715}uu-\tfrac{256}{15015},\\[5pt]&\qquad W^{\mathrm{VII}}=u^7-\tfrac{21}{13}u^5+\tfrac{105}{143}u^3-\tfrac{35}{429}u\ \text{etc.}\end{aligned}
\]
\SourceLeaf{196}
With a little attention, it becomes clear that each of $V,V',V'',V'''$, etc., $W,W',W'',W'''$, etc., is a polynomial in the indeterminate $u$; that the highest term in $V^{(m)}$ is $u^{m-1}$, and the powers $u^{m-2},u^{m-4},u^{m-6}$, etc., are absent; whereas the highest term in $W^{(m)}$ is $u^m$, and the powers $u^{m-1},u^{m-3},u^{m-5}$, etc., are absent. By what we have proved above, however, we shall have
\[
\begin{aligned}&\varphi=\frac1{WW'}+\frac1{3W'W''}+\frac{2\mathbin{.}2}{3\mathbin{.}3\mathbin{.}5W''W'''}\\[5pt]&+\frac{2\mathbin{.}2\mathbin{.}3\mathbin{.}3}{3\mathbin{.}3\mathbin{.}5\mathbin{.}5\mathbin{.}7W'''W''''}+\frac{2\mathbin{.}2\mathbin{.}3\mathbin{.}3\mathbin{.}4\mathbin{.}4}{3\mathbin{.}3\mathbin{.}5\mathbin{.}5\mathbin{.}7\mathbin{.}7\mathbin{.}9W''''W'''''}+\text{etc.}\end{aligned}
\]
and hence, in general,
\[
\begin{aligned}
\varphi-\frac{V^{(m)}}{W^{(m)}}&=\frac{2\mathbin{.}2\mathbin{.}3\mathbin{.}3\mathinner{.\,.\,.\,.\,.} m\mathbin{.}m}{3\mathbin{.}3\mathbin{.}5\mathbin{.}5\mathinner{.\,.\,.\,.}(2m-1)(2m+1)W^{(m)}W^{(m+1)}}\\[5pt]
&\quad+\frac{2\mathbin{.}2\mathbin{.}3\mathbin{.}3\mathinner{.\,.\,.\,.\,.}(m+1)(m+1)}{3\mathbin{.}3\mathbin{.}5\mathbin{.}5\mathinner{.\,.\,.\,.\,.}(2m+1)(2m+3)W^{(m+1)}W^{(m+2)}}\\[5pt]
&\quad+\text{etc.}
\end{aligned}
\]
If, therefore, $\varphi-\frac{V^{(m)}}{W^{(m)}}$ is expanded into a descending series, its first term will be
\[
=\frac{2\mathbin{.}2\mathbin{.}3\mathbin{.}3\mathinner{.\,.\,.\,.} m\mathbin{.}m u^{-(2m+1)}}{3\mathbin{.}3\mathbin{.}5\mathbin{.}5\mathinner{.\,.\,.\,.}(2m-1)(2m+1)}
\]
The product $\varphi W^{(m)}$, however, will consist of the polynomial $V^{(m)}$ and an infinite series whose first term is
\[
=\frac{2\mathbin{.}2\mathbin{.}3\mathbin{.}3\mathinner{.\,.\,.\,.\,.} mmu^{-(m+1)}}{3\mathbin{.}3\mathbin{.}5\mathbin{.}5\mathinner{.\,.\,.\,.\,.}(2m-1)(2m+1)}
\]
Thus the function $U$ of degree $n+1$, satisfying the condition established in the preceding article, has been found automatically: namely, that the product $\varphi U$ is free of the powers $u^{-1},u^{-2},u^{-3}\mathinner{.\,.\,.\,.\,.} u^{-(n+1)}$. It is precisely $W^{(n+1)}$; at the same time, it is clear that $U'$ is equal to $V^{(m+1)}$, and also that the first term of $U''$ is
\[
=\frac{2\mathbin{.}2\mathbin{.}3\mathbin{.}3\mathinner{.\,.\,.\,.\,.}(n+1)(n+1)}{3\mathbin{.}3\mathbin{.}5\mathbin{.}5\mathinner{.\,.\,.\,.\,.}(2n+1)(2n+3)}\mathbin{.}u^{-(n+2)}
\]
If, therefore, for $b,b',b''\mathinner{.\,.\,.\,.\,.} b^{(n)}$ we take the roots of the equation $W^{(n+1)}=0$, and the values of the coefficients $R,R',R''\mathinner{.\,.\,.\,.\,.} R^{(n)}$ are found by the rules given above, our formula for the integral will have accuracy extending to order $2n+1$, and its correction will be expressed approximately by
\[
\begin{aligned}&\frac1{2^{2n+2}}\mathbin{.}\frac{2\mathbin{.}2\mathbin{.}3\mathbin{.}3\mathinner{.\,.\,.\,.\,.}(n+1)(n+1)}{3\mathbin{.}3\mathbin{.}5\mathbin{.}5\mathinner{.\,.\,.\,.\,.}(2n+1)(2n+3)}L^{(2n+2)}\\[5pt]&=\frac{1\mathbin{.}1\mathbin{.}2\mathbin{.}2\mathbin{.}3\mathbin{.}3\mathinner{.\,.\,.\,.\,.}(n+1)(n+1)}{2\mathbin{.}6\mathbin{.}6\mathbin{.}10\mathbin{.}10\mathbin{.}14\mathinner{.\,.\,.\,.\,.}(4n+2)(4n+6)}L^{(2n+2)}\end{aligned}
\]
\SourceLeaf{197}
